\section{Loads and Boundary Conditions}
\label{sec:loads-bcs}

Loads and boundary conditions define how the finite-element model is forced and supported. FEMaster supports reusable named load and support collectors, step-local native load definitions, and persistent Abaqus-style condition histories. Named collectors are activated with \cmd{LOADS} and \cmd{SUPPORTS}. Unnamed native vector loads and native supports participate in the current analysis, while Abaqus-style \cmd{BOUNDARY}, single-DOF \cmd{CLOAD}, labeled \cmd{DLOAD}, and \cmd{DSLOAD} definitions update persistent conditions that are inherited by subsequent analyses.

For persistent direct conditions, \key{OP=MOD} (the default)
replaces matching components or targets; \key{OP=NEW} retires the
inherited entries in that condition family before applying the new
ones. For step-local native definitions, the supplied values apply
to the current analysis; \key{OP=NEW} also clears any inherited direct
entries of that family. Named collector definitions are reusable
between analyses. Thermal prescriptions maintain separate temperature,
heat-flux, and convection histories.

Throughout this section, remember that a numerical zero is a real physical value. For a native \cmd{SUPPORT}, \val{0} means that the corresponding displacement or rotation is constrained to zero, while an omitted or \val{NAN} component means that the degree of freedom remains unconstrained.

\keywordpage{AMPLITUDE}

\cmd{AMPLITUDE} defines a reusable scalar function that can scale a load as the analysis variable changes. In a transient analysis this variable is physical time. Both input readers accept up to four \((t,a)\) pairs per data line. \key{TYPE} selects \val{LINEAR}, \val{STEP}, or \val{NEAREST} interpolation; the default is \val{LINEAR}.

For linear interpolation, a value between two samples \((t_i,a_i)\) and \((t_{i+1},a_{i+1})\) is obtained from
\[
a(t)=a_i+\frac{t-t_i}{t_{i+1}-t_i}\left(a_{i+1}-a_i\right).
\]
\val{STEP} holds the preceding sample value between tabulated points. \val{NEAREST} uses the closest sample. The amplitude is dimensionless unless the user deliberately assigns another interpretation; it multiplies the referenced load magnitude.

An amplitude does not create a load by itself. It must be referenced by a supported load or boundary-condition keyword. This allows one time history to be reused by several load components.

The supported Abaqus form is \cmd{AMPLITUDE, DEFINITION=TABULAR}. \key{TIME=STEPTIME} and \key{VALUE=RELATIVE} are accepted by both readers. These options use the same tabular amplitude representation as the FEMaster interpolation modes.

In a transient analysis, amplitude sample times use the same time unit as \cmd{TIME}. A model expressed in seconds therefore interprets the amplitude abscissa in seconds.

\begin{keywordparameters}
\key{NAME} & required & Unique amplitude name. \\
\key{TYPE} & optional & \val{LINEAR}, \val{STEP}, or \val{NEAREST}; default \val{LINEAR}. \\
\key{DEFINITION} & optional & Supported value \val{TABULAR}. \\
\key{TIME} & optional & Supported value \val{STEPTIME}. \\
\key{VALUE} & optional & Supported value \val{RELATIVE}. \\
\end{keywordparameters}

\begin{keyworddata}
Both readers & One to four \val{time, value} pairs per physical line. \\
\end{keyworddata}

The first table describes a ramp from zero to full magnitude during the first 0.1 units of step time, followed by a constant factor. A load references the amplitude by its NAME.

\exampleheading

\begin{inputdeck}
*AMPLITUDE, NAME=RAMP,
 DEFINITION=TABULAR, TIME=STEPTIME
0.0, 0.0, 0.1, 1.0,
1.0, 1.0
\end{inputdeck}

\paragraph*{Stepwise interpolation}
\begin{inputdeck}
*AMPLITUDE, NAME=HOLD, TYPE=STEP
0., 0., 0.5, 1., 1., 1.
\end{inputdeck}

\keywordpage{SUPPORT}

Native \cmd{SUPPORT} prescribes generalized nodal displacements and rotations. Every row contains a target node or node set followed by six components in the order
\[
(U_x,U_y,U_z,R_x,R_y,R_z).
\]
A finite number prescribes that component. \val{0} therefore creates a homogeneous constraint. An omitted entry or \val{NAN} leaves the component free.

\key{SUPPORT_COLLECTOR} selects a reusable support collector activated by \cmd{SUPPORTS}. Within an active analysis the same six-component syntax can prescribe an unnamed step-local support. \key{OP=NEW} also clears inherited direct supports before these step-local values are applied.

\key{ORIENTATION} optionally specifies a local coordinate system. The prescribed components are then expressed in that basis, including cylindrical supports, sliders, or boundaries whose constrained direction is not aligned with the global axes.

Every constrained degree of freedom removes an admissible motion. Constraints without a corresponding physical support increase the model stiffness, while missing support constraints can leave a structural nullspace.

\begin{keywordparameters}
\key{SUPPORT_COLLECTOR}/\key{NAME} & optional in analysis & Named reusable support collector; required outside an analysis. \\
\key{OP} & optional & \val{MOD} (default) for ordinary selection; \val{NEW} also clears inherited direct supports. \\
\key{ORIENTATION} & optional & Named local basis for the prescribed components. \\
\end{keywordparameters}

\begin{keyworddata}
\val{target} & Node ID or node set. \\
Six values & \(U_x,U_y,U_z,R_x,R_y,R_z\); omitted or \val{NAN} means free. \\
\end{keyworddata}

This collector stores generalized nodal prescriptions for later activation by SUPPORTS. NAN means that a component is left free; zero fixes that component.

\exampleheading

\begin{inputdeck}
*SUPPORT, SUPPORT_COLLECTOR=BC
FIXED, 0., 0., 0., 0., 0., 0.
GUIDE, NAN, 0., 0., NAN, NAN, NAN
\end{inputdeck}

\paragraph*{Equivalent boundary conditions}
\begin{inputdeck}
*BOUNDARY
FIXED, 1, 6, 0.
GUIDE, 2, 3, 0.
\end{inputdeck}

\keywordpage{BOUNDARY}

\cmd{BOUNDARY} is the supported Abaqus displacement/rotation constraint keyword. \key{TYPE} currently supports \val{DISPLACEMENT} and defaults to that form.

Each data row contains a target, a first degree of freedom, an optional last degree of freedom, and an optional prescribed value. Abaqus structural DOF numbers are one based:
\[
1\rightarrow U_x,\quad
2\rightarrow U_y,\quad
3\rightarrow U_z,\quad
4\rightarrow R_x,\quad
5\rightarrow R_y,\quad
6\rightarrow R_z.
\]
If the last DOF is omitted, only the first is constrained. If the value is omitted it defaults to zero.

If the target nodes have an Abaqus \cmd{TRANSFORM}, the boundary components are interpreted in that local nodal basis. Nonzero prescribed values are currently supported for the documented static procedures. Amplitude-dependent prescribed values are subject to the restrictions described for the corresponding analysis procedure.

Native \cmd{SUPPORT} presents all six generalized components in one row. Abaqus \cmd{BOUNDARY} uses DOF ranges, so the same constraint can require several rows.

\begin{keywordparameters}
\key{TYPE} & optional & Supported form \val{DISPLACEMENT}; default \val{DISPLACEMENT}. \\
\key{AMPLITUDE} & optional & Named amplitude for the prescribed boundary history. \\
\key{OP} & optional & \val{MOD} (default) or \val{NEW} for direct support history. \\
\end{keywordparameters}

\begin{keyworddata}
\val{target, first_dof, last_dof, value} & DOF range 1--6; last DOF defaults to first and value defaults to zero. \\
\end{keyworddata}

The boundary condition data use numbered DOFs, with inclusive ranges. In this example one displacement is prescribed and two other components are fixed.

\exampleheading

\begin{inputdeck}
*BOUNDARY, TYPE=DISPLACEMENT
TIP, 1, 1, 1.5
TIP, 2, 3, 0.
\end{inputdeck}

\paragraph*{Support vector}
\begin{inputdeck}
*SUPPORT, SUPPORT_COLLECTOR=BC
TIP, 1.5, 0., 0., NAN, NAN, NAN
\end{inputdeck}

\keywordpage{CLOAD}

\cmd{CLOAD} applies concentrated nodal force and moment components. Native FEMaster writes all six generalized load components in one row:
\[
(F_x,F_y,F_z,M_x,M_y,M_z).
\]
Missing components default to zero.

The target may be an individual node or a node set. When a node set is used, the specified vector is applied to \emph{every} node in the set. It is not divided by the number of nodes. A reference node with a structural/distributing \cmd{COUPLING} instead transfers one definite resultant force or moment to a surface.

\key{ORIENTATION} optionally defines a local load basis. \key{AMPLITUDE} optionally multiplies the load by a named amplitude during analyses that evaluate time-dependent loading.

Abaqus \cmd{CLOAD} uses one DOF and one magnitude per row. DOFs 1--3 represent forces and 4--6 represent nodal moments. The supported input accepts real non-follower loading; imaginary and follower concentrated loads are not supported.

The native six-component vector form supplies a step-local load when no collector is named. Abaqus-style \val{target, dof, magnitude} rows instead update persistent load history per target and DOF. \key{OP=MOD} replaces matching DOF entries; \key{OP=NEW} retires inherited direct CLOAD entries before the current command.

When a concentrated nodal load represents a physically distributed load, the idealization can create local stress singularities or mesh-dependent peaks. Surface loading and distributing couplings represent load introduction over a finite area.

\begin{keywordparameters}
\key{NAME}/\key{LOAD_COLLECTOR} & optional in analysis & Named reusable collector; required outside an analysis. \\
\key{OP} & optional & \val{MOD} (default) or \val{NEW} for direct load history. \\
\key{ORIENTATION} & native optional & Named local load basis. \\
\key{AMPLITUDE} & optional & Named amplitude where supported. \\
\key{FOLLOWER} & Abaqus flag & Currently unsupported. \\
\key{REAL}/\key{IMAGINARY} & Abaqus flags & Real loading is supported; imaginary loading is not. \\
\end{keywordparameters}

A concentrated load can be expressed by DOF number or by a generalized force vector. When a node set is targeted, the specified force is applied separately to every node in the set.

\exampleheading

\begin{inputdeck}
*CLOAD
TIP, 1, 1000.
TIP, 2, -250.
\end{inputdeck}

\paragraph*{Vector notation}
\begin{inputdeck}
*CLOAD, LOAD_COLLECTOR=TIP_LOAD
TIP, 1000., -250., 0., 0., 0., 0.
\end{inputdeck}

\keywordpage{DLOAD}

\cmd{DLOAD} supports several distributed-load data forms in the
unified input language. The Abaqus-compatible \val{GRAV} label applies
mass-proportional acceleration, \val{BX}/\val{BY}/\val{BZ} prescribe
global body-force densities, and \val{P1}, \val{P2}, etc. apply
pressure to the numbered element face. FEMaster also accepts a
direct three-component surface-traction vector on a named surface.

The native \val{surface, tx, ty, tz} form supplies an unnamed step-local traction, while labeled Abaqus rows \val{element_region, label, magnitude, ...} create matching persistent history entries. \key{OP=MOD} updates those labeled entries by target and load label; \key{OP=NEW} clears inherited direct DLOAD entries.

The vector form specifies traction force per unit area in the selected
basis. Where \key{ORIENTATION} is given, those vector components are
resolved in the local coordinate system. \key{AMPLITUDE} optionally
scales a load in time.

The consistent nodal force generated by a traction \(\boldsymbol t\) is obtained conceptually from the surface integral
\[
\boldsymbol f_e
=\int_{A_e} N^T\boldsymbol t\,\mathrm dA,
\]
so mesh refinement changes the number of nodal force contributions but not the intended continuous traction magnitude.

For the \val{GRAV} form, the acceleration is
\[
\boldsymbol a=g\,\boldsymbol d.
\]
Selected regular \val{MASS} elements also receive concentrated
gravity loads. Auxiliary \cmd{POINTMASS} objects are not members
of an ELSET and therefore are not selected by \cmd{DLOAD, GRAV}.
For surface pressure on a named surface set, the corresponding
Abaqus-compatible command is \cmd{DSLOAD, P}.

Traction has units of force per unit area. Its total resultant equals the surface integral and therefore coincides numerically with the traction magnitude only when the loaded area is one in the chosen unit system.

\begin{keywordparameters}
\key{NAME}/\key{LOAD_COLLECTOR} & optional in analysis & Named reusable collector; required outside an analysis. \\
\key{OP} & optional & \val{MOD} (default) or \val{NEW} for direct load history. \\
\key{ORIENTATION} & native optional & Local traction basis. \\
\key{AMPLITUDE} & optional & Named time scaling where supported. \\
\key{REAL}/\key{IMAGINARY} & optional flags & Imaginary distributed loading is unsupported. \\
\end{keywordparameters}

\paragraph*{Gravity acceleration}

The four examples below describe distinct physical load types: gravity acceleration, force density, surface traction and element-face pressure. Their units and target regions are not interchangeable.

\begin{inputdeck}
*DLOAD
EALL, GRAV, 9810., 0., 0., -1.
\end{inputdeck}

\paragraph*{Body-force density}
\begin{inputdeck}
*DLOAD
EALL, BZ, -0.01
\end{inputdeck}

\paragraph*{Surface traction vector}
\begin{inputdeck}
*DLOAD
SIDE, 0., 2.5, 0.
\end{inputdeck}

\paragraph*{Element-face pressure}
\begin{inputdeck}
*DLOAD
SOLIDS, P2, 5.0
\end{inputdeck}

\keywordpage{PLOAD}

Native \cmd{PLOAD} applies scalar pressure to a surface or surface set. Unlike a vector traction, pressure acts normal to the selected finite-element surface. The sign therefore depends on the element side and surface orientation.

For pressure \(p\) and oriented unit normal \(\boldsymbol n\), the surface traction is conceptually
\[
\boldsymbol t=p\boldsymbol n.
\]
The resulting nodal forces are obtained by surface integration. The load is therefore independent of how many nodes happen to lie on the surface, provided the discretization converges to the same geometry.

\key{AMPLITUDE} may scale the pressure where supported. Abaqus input expresses the corresponding supported pressure through \cmd{DSLOAD} with type \val{P}.

The surface side and element connectivity determine the pressure direction. Reversing the load sign can compensate numerically for an opposite orientation, but the same orientation continues to govern contact and shell-result conventions.

\begin{keywordparameters}
\key{NAME}/\key{LOAD_COLLECTOR} & optional in analysis & Named reusable collector; required outside an analysis. \\
\key{OP} & optional & \val{MOD} (default) for ordinary selection; \val{NEW} also clears inherited direct pload entries. \\
\key{AMPLITUDE} & optional & Named pressure scaling function. \\
\end{keywordparameters}

The named surface must already exist. PLOAD applies a scalar pressure along its surface normal, whereas DSLOAD expresses the supported pressure through the P load label.

\exampleheading

\begin{inputdeck}
*PLOAD, LOAD_COLLECTOR=PRESSURE
INNER_FACE, 2.0
\end{inputdeck}

\paragraph*{Equivalent DSLOAD pressure}
\begin{inputdeck}
*DSLOAD
INNER_FACE, P, 2.0
\end{inputdeck}

\keywordpage{DSLOAD}

\cmd{DSLOAD} applies an Abaqus-compatible scalar pressure to a
named surface set. The supported load label is \val{P}; each data
record specifies a surface, the label \val{P}, and a pressure
magnitude. The surface normal determines the load direction.

The current implementation accepts only \key{FOLLOWER=NO};
follower pressure is not supported in nonlinear analyses.
\key{AMPLITUDE} may provide time scaling. Vector surface traction
(\val{TRVEC}) is not supported by \cmd{DSLOAD}; use the FEMaster
three-component \cmd{DLOAD} form instead.

\begin{keywordparameters}
\key{NAME}/\key{LOAD_COLLECTOR} & optional in analysis & Named reusable collector; required outside an analysis. \\
\key{AMPLITUDE} & optional & Named amplitude for time scaling. \\
\key{FOLLOWER} & optional & \val{NO}; follower surface loads are not supported. \\
\key{OP} & optional & \val{MOD} or \val{NEW} for load history. \\
\end{keywordparameters}

\begin{keyworddata}
Pressure & \val{surface, P, magnitude}. \\
\end{keyworddata}

The named \val{FACE_A} surface is an existing geometric region. The pressure magnitude is applied in the normal direction of that surface; its orientation therefore matters.

\exampleheading

\begin{inputdeck}
*DSLOAD, FOLLOWER=NO
FACE_A, P, 5.0
\end{inputdeck}


\paragraph*{Equivalent PLOAD pressure}
\begin{inputdeck}
*PLOAD
FACE_A, 5.0
\end{inputdeck}



\keywordpage{VLOAD}

Native \cmd{VLOAD} prescribes a three-component \emph{body-force density} on an element or element set. The components have units of force per unit volume (for example, \(\mathrm{N/mm^3}\)), \emph{not} acceleration. FEMaster integrates the load over the physical element volume without multiplying by material density:
\[
\boldsymbol f_e=\int_{\Omega_e}\boldsymbol N^T\boldsymbol b\,\mathrm d\Omega.
\]
The resulting force is independent of the mass density assigned to the material. This applies equally to globally defined loads and loads expressed in a local basis via \key{ORIENTATION}. \key{AMPLITUDE} optionally scales the prescribed force density.

For a gravity load, use a mass-proportional inertial load rather than entering acceleration components as \cmd{VLOAD}. The supported Abaqus \cmd{DLOAD} with \val{GRAV} interprets its magnitude and direction as acceleration and includes material density and selected regular concentrated masses via the \texttt{InertialLoad} path. The two commands therefore \emph{must not} be treated as numerically interchangeable. Auxiliary native \cmd{POINTMASS} definitions are not selected through a regular element set.

\begin{keywordparameters}
\key{NAME}/\key{LOAD\_COLLECTOR} & optional in analysis & Named reusable collector; required outside an analysis. \\
\key{OP} & optional & \val{MOD} (default) for ordinary selection; \val{NEW} also clears inherited direct vload entries. \\
\key{ORIENTATION} & native optional & Named local vector basis for the force density. \\
\key{AMPLITUDE} & native optional & Named time-scaling function. \\
\end{keywordparameters}

\begin{keyworddata}
\val{target, bx, by, bz} & Element or element-set target followed by three body-force-density components. \\
\end{keyworddata}

The supplied components are forces per unit volume rather than accelerations. Consequently the equivalent DLOAD BX/BY/BZ form does not multiply the value by mass density.

\exampleheading

\begin{inputdeck}
*VLOAD, LOAD_COLLECTOR=BODY_FORCE
EALL, 0., 0., -0.01
\end{inputdeck}

\paragraph*{DLOAD body-force component}
\begin{inputdeck}
*DLOAD
EALL, BZ, -0.01
\end{inputdeck}

\keywordpage{INERTIALOAD}

The current native keyword is spelled \cmd{INERTIALOAD}. It defines rigid-body inertial loading over an element set. Every row specifies a reference centre, translational acceleration, angular velocity \(\boldsymbol\omega\), and angular acceleration \(\boldsymbol\alpha\).

For position \(\boldsymbol r\) measured from the reference centre, the rigid-body acceleration field contains the familiar contributions
\[
\boldsymbol a(\boldsymbol r)
=
\boldsymbol a_0
+\boldsymbol\alpha\times\boldsymbol r
+\boldsymbol\omega\times(\boldsymbol\omega\times\boldsymbol r),
\]
for the prescribed instantaneous translational, tangential, and centrifugal terms. The corresponding inertial forces depend on the model's mass distribution.

Regular \val{TYPE=MASS} point elements are part of normal element topology. If such an element belongs to the target ELSET, its concentrated translational mass contributes directly to the rigid-body inertia load. A point-mass section with rotary inertia also contributes the corresponding Euler and gyroscopic moments when present through the native point-property path.

\key{CONSIDER_POINT_MASSES} has a narrower meaning: it controls whether the auxiliary point elements generated by native \cmd{POINTMASS, NSET=...} are added in addition to the ordinary target ELSET. The default is false. It does not disable regular \val{TYPE=MASS} elements already selected by the target region.

Density, concentrated masses, geometry, angular velocity, and acceleration must all use one consistent unit system. Centrifugal acceleration scales with \(\omega^2 r\), so a unit error in angular velocity or length can become especially severe.

The spelling \cmd{INERTIALOAD} matches the currently accepted native keyword exactly.

\begin{keywordparameters}
\key{NAME}/\key{LOAD_COLLECTOR} & optional in analysis & Named reusable collector; required outside an analysis. \\
\key{OP} & optional & \val{MOD} (default) for ordinary selection; \val{NEW} also clears inherited direct inertial-load entries. \\
\key{CONSIDER_POINT_MASSES} & optional & Include auxiliary native \cmd{POINTMASS} objects; Boolean, default \val{0}. \\
\end{keywordparameters}

\begin{keyworddata}
\val{target} & Element set receiving the inertial load; regular MASS elements in this set participate automatically. \\
\val{cx,cy,cz} & Reference centre. \\
\val{ax,ay,az} & Translational acceleration of the centre. \\
\val{wx,wy,wz} & Angular velocity. \\
\val{alphax,alphay,alphaz} & Angular acceleration. \\
\end{keyworddata}

The example describes inertial loading by translational acceleration, rotation and angular acceleration. It is meant for the previously defined region EALL, with optional inclusion of auxiliary point masses.

\exampleheading

\begin{inputdeck}
*INERTIALOAD, LOAD_COLLECTOR=ROT_INERTIA,
 CONSIDER_POINT_MASSES=1
EALL,
0., 0., 0.,
0., 0., 0.,
0., 0., 10.,
0., 0., 2.
\end{inputdeck}

\keywordpage{TEMPERATURE}

\cmd{TEMPERATURE} prescribes the absolute temperature of a node or node
set. It maintains a temperature-condition history that is distinct from
the model's initial temperature Field set by
\cmd{INITIALCONDITION, TYPE=TEMPERATURE}. Each data row provides one
compiled node reference or node-set name and one temperature value.
\key{OP=MOD} replaces matching targets and \key{OP=NEW} begins a
new set of temperature prescriptions.

\begin{keywordparameters}
\key{OP} & optional & \val{MOD} (default) or \val{NEW}. \\
\end{keywordparameters}
\begin{keyworddata}
\val{target, temperature} & Node ID or node set and absolute temperature. \\
\end{keyworddata}

For a model with the assembled node set \val{HEATED_NODES}, the command
prescribes a uniform temperature of 80 using the model's temperature units.

\exampleheading
\begin{inputdeck}
*TEMPERATURE, OP=MOD
HEATED_NODES, 80.
\end{inputdeck}

\keywordpage{HEATFLUX}

\cmd{HEATFLUX} prescribes a heat-input density on an existing surface
region. A positive heat flux \(q\) acts into the thermal balance and
contributes through \(\int_A N^T q\,\mathrm dA\) over the reference
surface. \key{AMPLITUDE} references a scalar history multiplying the
heat-input density. \key{OP=MOD} replaces the matching surface target;
\key{OP=NEW} starts a new heat-flux history.

\begin{keywordparameters}
\key{OP} & optional & \val{MOD} (default) or \val{NEW}. \\
\key{AMPLITUDE} & optional & Named amplitude multiplying the heat-input density. \\
\end{keywordparameters}
\begin{keyworddata}
\val{surface, heat_flux} & Named surface or surface set and heat input per unit area. \\
\end{keyworddata}

This model-level definition scales a heat flux on the existing
\val{HOT_FACE} surface from zero to its nominal value during one
unit of step progression.

\exampleheading
\begin{inputdeck}
*AMPLITUDE, NAME=POWER_RAMP,
 DEFINITION=TABULAR
0., 0., 1., 1.
*HEATFLUX, AMPLITUDE=POWER_RAMP
HOT_FACE, 0.25
\end{inputdeck}

\keywordpage{CONVECTION}

\cmd{CONVECTION} defines Newton cooling on an existing surface
region. The surface heat input is
\[
q=h(T_\infty-T),
\]
where \(h\) is the film coefficient, \(T_\infty\) is the absolute
ambient temperature, and \(T\) is the local surface temperature.
Each row supplies the target surface, film coefficient, and ambient
temperature. A named \key{AMPLITUDE} scales \(h\) while \(T_\infty\)
remains the specified value. \key{OP=MOD} replaces the matching
surface target; \key{OP=NEW} starts a new convection history.

\begin{keywordparameters}
\key{OP} & optional & \val{MOD} (default) or \val{NEW}. \\
\key{AMPLITUDE} & optional & Named amplitude multiplying the film coefficient. \\
\end{keywordparameters}
\begin{keyworddata}
\val{surface, h, ambient} & Named surface, film coefficient, and absolute ambient temperature. \\
\end{keyworddata}

The existing surface \val{COOLED_FACE} exchanges heat with an
environment at temperature 20 through a film coefficient of 0.015,
expressed in the model's consistent thermal units.

\exampleheading
\begin{inputdeck}
*CONVECTION
COOLED_FACE, 0.015, 20.
\end{inputdeck}
